\chapter{Kurzfassung}
Große Sprachmodelle (LLMs) sind weitgehend intransparent, was die Identifizierung jener Trainingsbeispiele erschwert, die eine bestimmte Vorhersage beeinflussen. Instanzbasierte Erklärungsmethoden, die auf Modellgradienten basieren, stellen aufgrund der hohen Dimensionalität moderner LLMs eine große rechnerische Herausforderung dar. Diese Arbeit untersucht, ob niedrigdimensionale Gradienten-Surrogate, die entweder durch die \textbf{Auswahl} (Selektion) einer Teilmenge von Schichtkomponenten oder durch die \textbf{Projektion} des gesamten Gradienten erstellt werden, instanzbasierte Erklärungen effektiv reproduzieren können. Die zentrale Fragestellung lautet: Ist es effektiver, auszuwählen oder zu projizieren?
\\\\
Zur Evaluierung instanzbasierter Erklärungen wird ein neuartiges, effizientes retrieval-basiertes Benchmark vorgestellt. Die Aufgabe besteht darin, ein originales Trainingsbeispiel anhand seiner paraphrasierten Version zu identifizieren, indem die Kosinus-Ähnlichkeit der jeweiligen Kostengradienten gemessen wird. Dieses Benchmark wird verwendet, um zwei Strategien zur Erstellung niedrigdimensionaler Gradientenrepräsentationen auf einem 1,2B-Parameter-LLM zu vergleichen: (1) ein gieriger Vorwärtsauswahl-Algorithmus, der eine kleine, architekturbasiert informierte Teilmenge einflussreicher Layerkomponenten identifiziert, und (2) die zufällige Projektion, die den vollständigen Gradienten auf einen dichten, niedrigdimensionalen Raum abbildet.
\\\\
Die Ergebnisse zeigen, dass für die Retrieval-Aufgabe eine gezielte Auswahl einer kleinen Teilmenge von Gradientenkomponenten überlegen ist. Eine gierig ausgewählte Teilmenge, die oft weniger als 5~\% der Modellparameter umfasst und hauptsächlich aus MLP- und Attention-Query/Output-Blöcken besteht, übertrifft nicht nur zufällige Projektionen, sondern erzielt auch eine höhere Abrufgenauigkeit als der vollständige Modellgradient selbst. Dies deutet darauf hin, dass der vollständige Gradient für diese Aufgabe ein verrauschtes und suboptimales Signal sein kann. Während die zufällige Projektion die globale Geometrie des Gradientenraums effektiv bewahrt, ist sie für das spezifische Abrufziel rechenintensiver und ungenauer.
\\\\
Die Resultate liefern eine klare Antwort: Für die Entwicklung effizienter und effektiver instanzbasierter Erklärungswerkzeuge, die als Retrieval-Aufgabe konzipiert sind, ist die Auswahl einer kleinen, architekturbewussten Repräsentation des Gradienten genauer und recheneffizienter als die Projektion des gesamten Gradienten. Diese Arbeit leistet einen Beitrag durch ein effizientes Benchmark für instanzbasierte Erklärungen, einen prinzipiengeleiteten Vergleich von Dimensionsreduktionsstrategien und empirische Belege, die eine gezielte Auswahl gegenüber einer holistischen Projektion befürworten.
\clearpage